\begin{operation}\label{operation_BIU}
We note the operation given in the proof of \cite[Theorem~1.1]{BIU}, which we use in (zig-5) and (zag-5).
\begin{itemize}\itemsep=0pt
\item[\rm (a)] The dimer model $\overline{\nu}^\zig_\calX(\Gamma, \{z_1,\dots,z_r\})$ (resp.\ $\overline{\nu}^\zag_\calY(\Gamma, \{z_1,\dots,z_r\})$), which is obtained by
applying the operations (zig-1)--(zig-4) (resp.~(zag-1)--(zag-4)) to the reduced consistent dimer model $\Gamma$, sometimes contains zigzag paths having self-intersections in the universal cover.
In this case, we use the operation given in the proof of \cite[Theorem~1.1]{BIU}; that is, we remove all edges at self-intersections (see Figure~\ref{fig_remove_self}).
We note that this operation does not change the slope of the argued zigzag path.
\begin{figure}[h!]\centering
\scalebox{0.5}{
\begin{tikzpicture}
\newcommand{\edgewidth}{0.05cm} %line_width_of _edges
\newcommand{\nodewidth}{0.05cm} %line_width_around_nodes
\newcommand{\noderad}{0.16} %radius_of_node

%coordinate setting
\coordinate (B1) at (0,0); \coordinate (B2) at (2,3);
\coordinate (W1) at (0,2); \coordinate (W2) at (2,-1);
\path (W1) ++(135:1.2cm) coordinate (W1a); \path (B1) ++(225:1.2cm) coordinate (B1a);
\path (W2) ++(45:1.2cm) coordinate (W2a); \path (W2) ++(0:1.2cm) coordinate (W2b);
\path (B2) ++(0:1.2cm) coordinate (B2a); \path (B2) ++(315:1.2cm) coordinate (B2b);

\node (original) at (0,0) {
\begin{tikzpicture}
%edge
\draw [line width=\edgewidth] (B1)--(B1a); \draw [line width=\edgewidth] (W1)--(W1a);
\draw [line width=\edgewidth] (B1)--(W1); \draw [line width=\edgewidth] (B1)--(W2); \draw [line width=\edgewidth] (B2)--(W1);
\draw [line width=\edgewidth] (W2)--(W2a); \draw [line width=\edgewidth] (W2)--(W2b);
\draw [line width=\edgewidth] (B2)--(B2a); \draw [line width=\edgewidth] (B2)--(B2b);
%white node
\draw [line width=\nodewidth, fill=white] (W1) circle [radius=\noderad] ; \draw [line width=\nodewidth, fill=white] (W2) circle [radius=\noderad] ;
%black node
\draw [line width=\nodewidth, fill=black] (B1) circle [radius=\noderad] ; \draw [line width=\nodewidth, fill=black] (B2) circle [radius=\noderad] ;
%zigzag
\path (B1) ++(135:0.25cm) coordinate (B1+); \path (B1+) ++(225:1.3cm) coordinate (B1++);
\path (W1) ++(225:0.25cm) coordinate (W1+); \path (W1+) ++(135:1.3cm) coordinate (W1++);
\path (B1) ++(45:0.25cm) coordinate (B1-);
\path (W1) ++(315:0.25cm) coordinate (W1-);
\path (B2) ++(270:0.25cm) coordinate (B2+); \path (B2+) ++(0:1.6cm) coordinate (B2++);
\path (W2) ++(90:0.25cm) coordinate (W2+); \path (W2+) ++(0:1.6cm) coordinate (W2++);
\draw [->, rounded corners, line width=0.08cm, red] (B1++)--(B1+)--(W1-)--(B2+)--(B2++);
\draw [->, rounded corners, line width=0.08cm, red] (W1++)--(W1+)--(B1-)--(W2+)--(W2++);
\end{tikzpicture}};

\node (remove) at (10,0) {
\begin{tikzpicture}
%edge
\draw [line width=\edgewidth] (B1)--(B1a); \draw [line width=\edgewidth] (W1)--(W1a);
\draw [line width=\edgewidth] (B1)--(W2); \draw [line width=\edgewidth] (B2)--(W1);
\draw [line width=\edgewidth] (W2)--(W2a); \draw [line width=\edgewidth] (W2)--(W2b);
\draw [line width=\edgewidth] (B2)--(B2a); \draw [line width=\edgewidth] (B2)--(B2b);
%white node
\draw [line width=\nodewidth, fill=white] (W1) circle [radius=\noderad] ; \draw [line width=\nodewidth, fill=white] (W2) circle [radius=\noderad] ;
%black node
\draw [line width=\nodewidth, fill=black] (B1) circle [radius=\noderad] ; \draw [line width=\nodewidth, fill=black] (B2) circle [radius=\noderad] ;
%zigzag
\path (B2) ++(270:0.3cm) coordinate (B2+); \path (W2) ++(90:0.3cm) coordinate (W2+);
\path (B2+) ++(0:1.6cm) coordinate (B2++); \path (W2+) ++(0:1.6cm) coordinate (W2++);
\path (B1) ++(90:0.3cm) coordinate (B1n); \path (B1n) ++(225:1.5cm) coordinate (B1n+);
\path (W1) ++(270:0.3cm) coordinate (W1s); \path (W1s) ++(135:1.5cm) coordinate (W1s+);
\draw [->, rounded corners, line width=0.08cm, red] (B1n+)--(B1n)--(W2+)--(W2++);
\draw [->, rounded corners, line width=0.08cm, red] (W1s+)--(W1s)--(B2+)--(B2++);
\end{tikzpicture}};

\path (original) ++(0:4cm) coordinate (original+); \path (remove) ++(180:4cm) coordinate (remove+);
\draw [->, line width=\edgewidth] (original+)--(remove+);
\end{tikzpicture}
}
\caption{An example of removing a self-intersection of a zigzag path.}\label{fig_remove_self}
\end{figure}

After this process, there might be a connected component of the resulting bipartite graph that is contained in a simply-connected domain in $\TT$.
In that case, we remove such a~connected component.
We note that this removal does not affect our purpose, because our main concern is the PM polygon which is recovered from the slopes of zigzag paths,
and the slope of the zigzag path corresponding to the argued connected component is trivial.

\item[\rm (b)] On the other hand, the dimer model $\overline{\nu}^\zig_\calX(\Gamma, \{z_1,\dots,z_r\})$ (resp.\ $\overline{\nu}^\zag_\calY(\Gamma, \{z_1,\dots,z_r\})$)
might have a pair of zigzag paths on the universal cover that intersect with each other in the same direction more than once.
In this case, we use another operation given in the proof of \cite[Theorem~1.1]{BIU}, that is,
we choose any such pair of zigzag paths and remove pairs of consecutive intersections of this pair of zigzag paths (see Figure~\ref{fig_remove_pair}).
We note that this operation does not change the slopes of zigzag paths and the resulting bipartite graph is also a dimer model
because $\overline{\nu}^\zig_\calX(\Gamma, \{z_1,\dots,z_r\})$ (resp.\ $\overline{\nu}^\zag_\calY(\Gamma, \{z_1,\dots,z_r\})$) satisfies the strong marriage condition, as we will see in Proposition~\ref{prop_nondegenerate}.
\begin{figure}[h!]\centering
\scalebox{0.5}{
\begin{tikzpicture}
\newcommand{\edgewidth}{0.05cm} %line_width_of _edges
\newcommand{\nodewidth}{0.05cm} %line_width_around_nodes
\newcommand{\noderad}{0.16} %radius_of_node

%coordinate setting
\coordinate (B1) at (0,0); \coordinate (B2) at (2,2); \coordinate (B3) at (4,0);
\coordinate (W1) at (0,2); \coordinate (W2) at (2,0); \coordinate (W3) at (4,2);

\path (W1) ++(135:1.2cm) coordinate (W1a); \path (B1) ++(225:1.2cm) coordinate (B1a);
\path (W3) ++(45:1.2cm) coordinate (W3a); \path (B3) ++(315:1.2cm) coordinate (B3a);

\node (original) at (0,0) {
\begin{tikzpicture}
%edge
\draw [line width=\edgewidth] (B1)--(W1)--(B2)--(W3)--(B3)--(W2)--(B1); \draw [line width=\edgewidth] (B2)--(W2);
\draw [line width=\edgewidth] (B1)--(B1a); \draw [line width=\edgewidth] (W1)--(W1a);
\draw [line width=\edgewidth] (B3)--(B3a); \draw [line width=\edgewidth] (W3)--(W3a);
%white node
\draw [line width=\nodewidth, fill=white] (W1) circle [radius=\noderad] ; \draw [line width=\nodewidth, fill=white] (W2) circle [radius=\noderad] ;
\draw [line width=\nodewidth, fill=white] (W3) circle [radius=\noderad] ;
%black node
\draw [line width=\nodewidth, fill=black] (B1) circle [radius=\noderad] ; \draw [line width=\nodewidth, fill=black] (B2) circle [radius=\noderad] ;
\draw [line width=\nodewidth, fill=black] (B3) circle [radius=\noderad] ;

%zigzag
\path (W1) ++(45:0.25cm) coordinate (W1+); \path (W1+) ++(135:1.3cm) coordinate (W1++);
\path (W1) ++(225:0.25cm) coordinate (W1-);
\path (W3) ++(135:0.25cm) coordinate (W3+); \path (W3+) ++(45:1.3cm) coordinate (W3++);
\path (W3) ++(315:0.25cm) coordinate (W3-);
\path (B1) ++(45:0.25cm) coordinate (B1+);
\path (B1) ++(135:0.25cm) coordinate (B1-); \path (B1-) ++(225:1.3cm) coordinate (B1--);
\path (B3) ++(135:0.25cm) coordinate (B3+);
\path (B3) ++(45:0.25cm) coordinate (B3-); \path (B3-) ++(315:1.3cm) coordinate (B3--);
\draw [->, rounded corners, line width=0.08cm, blue] (B1--)--(B1-)--(W1-)--(W3-)--(B3-)--(B3--);
\draw [->, rounded corners, line width=0.08cm, red] (W1++)--(W1+)--(B1+)--(B3+)--(W3+)--(W3++);
\end{tikzpicture}
};

\node (remove) at (10,0) {
\begin{tikzpicture}
%edge
\draw [line width=\edgewidth] (W1)--(B2)--(W3); \draw [line width=\edgewidth] (B1)--(W2)--(B3);
\draw [line width=\edgewidth] (B2)--(W2);
\draw [line width=\edgewidth] (B1)--(B1a); \draw [line width=\edgewidth] (W1)--(W1a);
\draw [line width=\edgewidth] (B3)--(B3a); \draw [line width=\edgewidth] (W3)--(W3a);
%white node
\draw [line width=\nodewidth, fill=white] (W1) circle [radius=\noderad] ; \draw [line width=\nodewidth, fill=white] (W2) circle [radius=\noderad] ;
\draw [line width=\nodewidth, fill=white] (W3) circle [radius=\noderad] ;
%black node
\draw [line width=\nodewidth, fill=black] (B1) circle [radius=\noderad] ; \draw [line width=\nodewidth, fill=black] (B2) circle [radius=\noderad] ;
\draw [line width=\nodewidth, fill=black] (B3) circle [radius=\noderad] ;
%zigzag
\path (B1) ++(315:0.25cm) coordinate (B1-a); \path (B1-a) ++(225:1.3cm) coordinate (B1--a);
\path (B3) ++(225:0.25cm) coordinate (B3-a); \path (B3-a) ++(315:1.3cm) coordinate (B3--a);
\draw [->, rounded corners, line width=0.08cm, red] (W1++)--(W1+)--(W3+)--(W3++);
\draw [->, rounded corners, line width=0.08cm, blue] (B1--a)--(B1-a)--(B3-a)--(B3--a);
\end{tikzpicture}
};

\path (original) ++(0:4cm) coordinate (original+); \path (remove) ++(180:4cm) coordinate (remove+);
\draw [->, line width=\edgewidth] (original+)--(remove+);
\end{tikzpicture}}
\caption{An example of removing a pair of consecutive intersections of zigzag paths.}\label{fig_remove_pair}
\end{figure}
\end{itemize}
\end{operation}

Since the dimer model $\overline{\nu}^\zig_\calX(\Gamma, \{z_1,\dots,z_r\})$ $($resp.\ $\overline{\nu}^\zag_\calY(\Gamma, \{z_1,\dots,z_r\})$$)$ is non-degenerate by Proposition~\ref{prop_nondegenerate}
and since it does not contain a homologically trivial zigzag path (see the proofs of Propositions~\ref{deform_vector}, \ref{zigzag_afterdeform1} and~\ref{zigzag_afterdeform2}),
we can produce a dimer model satisfying the conditions in Definition~\ref{def_consistent}
from $\overline{\nu}^\zig_\calX(\Gamma, \{z_1,\dots,z_r\})$ (resp.\ $\overline{\nu}^\zag_\calY(\Gamma, \{z_1,\dots,z_r\})$)
by iterated application of Operation~\ref{operation_BIU}.
Thus, $\nu^\zig_\calX(\Gamma, \{z_1,\dots,z_r\})$ and $\nu^\zag_\calY(\Gamma, \{z_1,\dots,z_r\})$ are consistent dimer models, but those are not necessarily isoradial even if $\Gamma$ is isoradial (see Example~\ref{ex_not_isoradial}).
Furthermore, since the operation (join) does not change the slopes of zigzag paths, this proves the following proposition.

\begin{Proposition}\label{prop_make_consistent}
Let the notation be the same as Definitions {\rm \ref{def_exdeformation_zig}} and {\rm \ref{def_exdeformation_zag}}.
Then, we have the following.
\begin{itemize}\itemsep=0pt
\item[\rm (1)] The dimer models $\nu^\zig_\calX(\Gamma, \{z_1,\dots,z_r\})$ and $\nu^\zag_\calY(\Gamma, \{z_1,\dots,z_r\})$ are consistent.
\item[\rm (2)] The set of slopes of zigzag paths on $\nu^\zig_\calX(\Gamma, \{z_1,\dots,z_r\})$ $($resp.\ $\nu^\zag_\calY(\Gamma, \{z_1,\dots,z_r\})$$)$
is the same as that of $\overline{\nu}^\zig_\calX(\Gamma, \{z_1,\dots,z_r\})$ $($resp.\ $\overline{\nu}^\zag_\calY(\Gamma, \{z_1,\dots,z_r\})$$)$.
\end{itemize}
\end{Proposition}
